% Gabarit du livrable J1 -- Note de cadrage
% Remplacez chaque [À compléter] ; supprimez les lignes d'aide en gris avant le dépôt.
% Compiler avec XeLaTeX ou LuaLaTeX (Overleaf : Menu > Compilateur > XeLaTeX)
\documentclass[11pt,a4paper]{article}
\usepackage{fontspec}
\usepackage{polyglossia}
\setmainlanguage{french}
\usepackage[a4paper,margin=2cm,headheight=34pt,includehead]{geometry}
\usepackage{graphicx}
\usepackage{xcolor}
\usepackage{tabularx,booktabs,array}
\usepackage{enumitem}
\usepackage{titlesec}
\usepackage{fancyhdr}
\usepackage{xurl}
\usepackage[hidelinks]{hyperref}

\definecolor{bleu}{HTML}{0B5C8C}
\definecolor{encre}{HTML}{172131}
\definecolor{gris}{HTML}{5E6B7E}
\hypersetup{colorlinks=true,linkcolor=bleu,urlcolor=bleu}

\setlength{\parindent}{0pt}
\setlength{\parskip}{5pt}
\setlength{\emergencystretch}{3em}
\renewcommand{\arraystretch}{1.35}
\titleformat{\section}{\large\bfseries\sffamily\color{bleu}}{\thesection.}{0.5em}{}
\titleformat{\subsection}{\normalsize\bfseries\sffamily\color{encre}}{\thesubsection}{0.5em}{}
\newcolumntype{L}[1]{>{\raggedright\arraybackslash}p{#1}}
\newcolumntype{Y}{>{\raggedright\arraybackslash}X}

% ---- À renseigner une fois ----
\newcommand{\projet}{P[XX]}
\newcommand{\titreprojet}{[Titre du projet]}
\newcommand{\filiere}{[Big Data Analytics / Full-Stack]}
\newcommand{\groupe}{[Nom du groupe]}
\newcommand{\depot}{[https://github.com/...]}
% Ligne d'aide, à supprimer avant le dépôt :
\newcommand{\aide}[1]{{\small\color{gris}\itshape #1}\par}
\newcommand{\todo}{{\color{red!70!black}[À compléter]}}

\pagestyle{fancy}
\fancyhf{}
\fancyhead[L]{\includegraphics[height=24pt]{logo_uemf.png}\hspace{8pt}\includegraphics[height=24pt]{logo_eidia.png}}
\fancyhead[R]{\small\sffamily\color{gris}J1 · Note de cadrage\\\projet{} · \groupe}
\fancyfoot[R]{\small\sffamily\color{gris}\thepage}

\begin{document}

% ================= Page de garde =================
\begin{titlepage}
\centering
\includegraphics[height=60pt]{logo_uemf.png}\hspace{16pt}\includegraphics[height=60pt]{logo_eidia.png}\par
\vspace{2.5cm}
{\sffamily\large\color{gris}Module Gestion des bases de données · Projet de semestre 2026-2027\par}
\vspace{0.6cm}
{\sffamily\Huge\bfseries\color{bleu}Livrable J1 · Note de cadrage\par}
\vspace{0.4cm}
{\sffamily\LARGE\bfseries\color{encre}\projet{} · \titreprojet\par}
\vspace{0.4cm}
{\sffamily\large Filière \filiere\par}
\vspace{1.8cm}
\begin{tabularx}{0.6\linewidth}{Y}
\toprule
\textbf{Membres du groupe} \\
\midrule
\todo \\
\todo \\
\todo \\
\todo \\
\bottomrule
\end{tabularx}
\vfill
Dépôt Git : \url{\depot}\par
Date de dépôt : \todo
\end{titlepage}

% ================= 1 =================
\section{Présentation du sujet}
\aide{½ page. Qui est le commanditaire, qui utilisera le système, quelle décision il aidera à prendre, pourquoi c'est important dans la vie réelle.}
\todo

% ================= 2 =================
\section{Problématique et objectifs}
\subsection{Problématique}
\aide{Une seule phrase interrogative, reformulée par le groupe.}
\todo
\subsection{Périmètre}
\begin{tabularx}{\linewidth}{L{3.5cm} Y Y}
\toprule
& \textbf{Inclus} & \textbf{Exclu} \\
\midrule
Période & \todo & \todo \\
Zone & \todo & \todo \\
Population / objets & \todo & \todo \\
\bottomrule
\end{tabularx}
\subsection{Objectifs}
Objectif général : \todo

\begin{tabularx}{\linewidth}{L{0.6cm} Y L{6.5cm}}
\toprule
\textbf{N°} & \textbf{Objectif spécifique (verbe d'action)} & \textbf{Indicateur de réussite} \\
\midrule
O1 & \todo & \todo \\
O2 & \todo & \todo \\
O3 & \todo & \todo \\
\bottomrule
\end{tabularx}

% ================= 3 =================
\section{Questions métier}
\aide{Les 5 questions du cahier des charges + au moins une question ajoutée par le groupe (Q6).}
{\small
\begin{tabularx}{\linewidth}{L{0.6cm} Y L{4.2cm} L{4.2cm}}
\toprule
\textbf{N°} & \textbf{Question} & \textbf{Indicateur / calcul} & \textbf{Données nécessaires} \\
\midrule
Q1 & \todo & & \\
Q2 & & & \\
Q3 & & & \\
Q4 & & & \\
Q5 & & & \\
Q6 & & & \\
\bottomrule
\end{tabularx}}

% ================= 4 =================
\section{Les données}
\subsection{Inventaire des sources}
\aide{Au moins 4 sources de 3 types différents, dont une API et le kit de l'enseignant.}
{\small
\begin{tabularx}{\linewidth}{L{2.4cm} L{1.6cm} L{1.5cm} Y L{1.7cm} L{1.5cm} L{1.5cm} L{1.6cm}}
\toprule
\textbf{Source} & \textbf{Type} & \textbf{Format} & \textbf{Accès (lien, clé)} & \textbf{Licence} & \textbf{Période} & \textbf{Volume} & \textbf{Mise à jour} \\
\midrule
\todo & & & & & & & \\
 & & & & & & & \\
 & & & & & & & \\
 & & & & & & & \\
\bottomrule
\end{tabularx}}

\subsection{Description de chaque source}
\aide{Une fiche par source. Copiez le bloc ci-dessous autant de fois que nécessaire.}

\textbf{Source : \todo}\\
Rôle (questions concernées) : \todo\\
Structure observée (fichiers, onglets, routes, lignes, colonnes) : \todo\\
Clé de liaison avec les autres sources : \todo

{\small
\begin{tabularx}{\linewidth}{L{2.6cm} L{2cm} L{1.6cm} L{2.6cm} Y}
\toprule
\textbf{Champ} & \textbf{Type observé} & \textbf{Unité} & \textbf{Exemple} & \textbf{Signification} \\
\midrule
\todo & & & & \\
 & & & & \\
\bottomrule
\end{tabularx}}

\subsection{Choix et justification}
\aide{Pourquoi ces sources ; sources envisagées puis écartées et raison ; schéma de liaison entre les sources (figure).}
\todo

\subsection{Premier diagnostic de qualité}
\aide{Au moins 5 anomalies réellement observées sur vos échantillons.}
{\small
\begin{tabularx}{\linewidth}{L{2.6cm} L{4cm} L{3cm} Y}
\toprule
\textbf{Source} & \textbf{Anomalie observée} & \textbf{Exemple} & \textbf{Traitement envisagé} \\
\midrule
\todo & & & \\
 & & & \\
 & & & \\
 & & & \\
 & & & \\
\bottomrule
\end{tabularx}}

\subsection{Données sensibles et contraintes}
\aide{Données personnelles ou de santé, loi 09-08 (CNDP), RGPD le cas échéant, licences, limites d'appels des API, gestion des clés.}
\todo

% ================= 5 =================
\section{Méthodologie}
\subsection{Pipeline envisagé}
\aide{Schéma de votre pipeline (sources, raw, stg, ods, dw, restitution, sécurité) avec l'outil de chaque étape.}
\todo
\subsection{Outils et conventions}
\todo
% ================= 6 =================
\section{Références}
\aide{Organisme, titre, lien, date de consultation, licence.}
\begin{enumerate}
\item \todo
\end{enumerate}

\appendix
\section{Annexes}
\aide{Extrait de chaque fichier, réponse d'API brute (tronquée), captures.}
\todo

\end{document}
